\section{Verification}\label{sec:verification}


\cref{fact:H} asserts that an explicit triple $(\Tbl,\Lv,10)$ is a valid certificate. This
section specifies what has to be checked (\cref{ssec:what}), gives the two algorithms that do it in a form
precise enough to be implemented from the paper (\cref{ssec:algverify,ssec:algleaf}), reports the independent
implementations and their results (\cref{ssec:impl}), summarizes the audits of the main verifier
(\cref{ssec:audits}) and several cross-checks that are not part of the proof (\cref{ssec:cross}), and
finally states plainly what has to be trusted (\cref{ssec:trust}). The details of the audits and of the
cross-checks are given in \cref{sec:supp-checks}.

\subsection{What must be checked}\label{ssec:what}

Let $(\Tbl,\Lv,\Rmax)$ be a certificate in the sense of \cref{def:certificate}: $\Tbl$ is a map from a
finite set of canonical Q-states to cells and $\Lv$ a finite set of canonical P-states. Its \emph{reachable
set} $V$ is the smallest set of canonical states that contains the root $\Pst(\emptyset,\emptyset,0)$ (R0)
and is closed under two rules: if a P-state of $V$ is not in $\Lv$, the canonical forms of all its
non-\VOID{} children of types (a), (c), (d) are in $V$ (R1); if a Q-state $q=\Qst(B,F,t)$ of $V$ has a table
entry $m=\Tbl(q)\notin B\cup F$, the canonical forms of its child (i) and, if it exists and is not \VOID, of
its child (ii) with respect to $m$ are in $V$ (R2). The certificate is \emph{valid} if
\begin{enumerate}[label=(V\arabic*),leftmargin=3em]
\item $\LEAF(B,F,t)$ holds for every P-state $\Pst(B,F,t)\in V\cap\Lv$;
\item every P-state $\Pst(B,F,t)\in V\setminus\Lv$ has $t<\Rmax$, and $\LOST(B,F,t)$ is false;
\item for every Q-state $q=\Qst(B,F,t)\in V$ the entry $\Tbl(q)$ is defined and is a cell outside $B\cup F$.
\end{enumerate}
By \cref{lem:certificate}, every state in the reachable set of a valid certificate is \WON; in particular
the root is \WON.

Validity is a finite, purely combinatorial statement. A verifier has to enumerate $V$ and to evaluate
three predicates: admissibility (for the \VOID{} test), $\LEAF$ and $\LOST$. The two files are treated as
\emph{untrusted hints}: for every possible content of the files, acceptance must imply that the root is
\WON. Consequently a defect of a verifier is dangerous only if it makes the verifier (1)~omit a child,
(2)~declare a non-\VOID{} child \VOID, (3)~evaluate $\LEAF$ as true or (4)~$\LOST$ as false wrongly,
(5)~accept an illegal table move, (6)~identify two inequivalent states, (7)~leave a queued state
unprocessed, or (8)~report success without having done all of this. Errors in parsing and truncated files can only
turn a valid certificate into a rejected one. The same holds for a canonical form that is not minimal, as
long as it is the image of the state under some element of $\G$ (\cref{rem:any-image}); a routine that
returns anything else falls under (1) or (6).

\subsubsection*{Defects of the input files}
\Cref{alg:verify} takes a map $\Tbl$ and a set $\Lv$ as its input; a program reads them from two files
and has to decide what to do when a file does not describe such objects. We distinguish two aims.

\emph{Soundness} (acceptance implies that the root is \WON) requires no test of the files at all, as long
as the verifier uses at most one move for every Q-state that it processes, recomputes $\LEAF$ for every
state that it treats as a leaf, and rejects when the entry of a Q-state that it processes is missing or is
not a cell outside $B\cup F$. An entry that is malformed, that is not canonical or that belongs to an
unreachable state may be ignored: ignoring an entry can only remove a hint. If a key occurs twice with
different moves, the verifier may reject, or it may use one of the two moves; the move that is used is
checked like any other.

\emph{Establishing \cref{fact:H} as stated} requires more, because the Computed Fact speaks of the map
and of the set that the two files define and says that every entry is reached. For this a verifier must
reject, or report so that the user can reject, each of the following: an entry that is not well formed in
the sense of \cref{ssec:fmt-leaves}; an entry whose state is not canonical; a key that occurs more than once
(in the table: also with the same move); and an entry that is never reached. All four counts are zero
for the two certificate files, in every program of \cref{tab:verifiers} that reports them. The main
verifier accepts some of these defects and counts them (\cref{ssec:supp-issues}), which is why its counters
have to be read together with its verdict (\cref{ssec:impl}).

Verifiers work with canonical forms. This is sound because every notion involved is defined through the
incidence of cells and lines and is therefore invariant under the group $\G$
(\cref{lem:game-invariance,lem:canon}): a state is \WON{} if and only if its image under $g\in\G$ is, and
the children of the image are the images of the children. A verifier that stores an arbitrary image of a
state instead of the minimal one is still sound, although \cref{prop:verify-correct} does not apply to it
literally; see \cref{rem:any-image}.

\subsection{Expanding the tree}\label{ssec:algverify}

\Cref{alg:verify} is the verifier. We write ``\textbf{require}~$\langle$condition$\rangle$'' for ``if the
condition does not hold, stop and return \reject''.

\begin{algorithm}[t]
\DontPrintSemicolon
\SetKwProg{Proc}{procedure}{}{}
\SetKwFunction{Push}{Push}
\SetKw{Require}{require}
\SetKw{Continue}{continue}
\KwIn{a table $\Tbl$ from canonical Q-states to cells; a set $\Lv$ of canonical P-states; an integer $\Rmax$}
\KwOut{\accept{} or \reject}
$W\gets\{\canon\Pst(\emptyset,\emptyset,0)\}$;\quad $\mathit{Seen}\gets W$\;
\While{$W\neq\emptyset$}{
  remove some state $s$ from $W$\;
  \uIf{$s$ is a P-state $\Pst(B,F,t)$}{
    \If{$s\in\Lv$}{\Require $\LEAF(B,F,t)$ \tcp*[r]{\cref{alg:leaf}}
       \Continue\;}
    \Require $t<\Rmax$\;
    \Require $\LOST(B,F,t)$ is false \tcp*[r]{\cref{def:lost}}\nllabel{ln:lost}
    \Push{$\Qst(B,F,t+1)$} \tcp*[r]{child (a)}\nllabel{ln:childa}
    \ForEach{live line $l$ and cell $e\in l\setminus F$ with $|F\cup(l\setminus\{e\})|\le t+1$}{
       \Push{$\Pst(B\cup\{e\},\,F\cup(l\setminus\{e\}),\,t+1)$} \tcp*[r]{child (c)}
    }
    \ForEach{cell $c\notin B\cup F$ and unordered pair $\{a,b\}$ of distinct live lines through $c$}{
      \ForEach{$p\in a\setminus(F\cup\{c\})$ and $q\in b\setminus(F\cup\{c\})$}{
        $Q_a\gets a\setminus\{c,p\}$;\quad $Q_b\gets b\setminus\{c,q\}$\;
        \If{$|F\cup Q_a\cup Q_b|\le t+1$}{
          \Push{$\Pst(B\cup\{c\},\,F\cup Q_a\cup Q_b,\,t+1)$} \tcp*[r]{child (d)}\nllabel{ln:childd}
        }
      }
    }
  }
  \Else(\tcp*[h]{$s$ is a Q-state $\Qst(B,F,t)$}){
    \Require $\Tbl(s)$ is defined;\quad $m\gets\Tbl(s)$\;
    \Require $m$ is a cell and $m\notin B\cup F$\;
    \Push{$\Pst(B\cup\{m\},F,t)$} \tcp*[r]{child (i)}
    \lIf{$|F|<t$}{\Push{$\Qst(B,F\cup\{m\},t)$} \tcp*[f]{child (ii)}}
  }
}
\Return \accept\;
\BlankLine
\Proc{\Push{$s'$}}{
  \lIf{$s'$ is \VOID, i.e.\ its set $F'$ is not $B'$-admissible}{\Return \tcp*[f]{never for a child (a) or (i)}}\nllabel{ln:void}
  $s''\gets\canon(s')$\;
  \lIf{$s''\notin\mathit{Seen}$}{add $s''$ to $\mathit{Seen}$ and to $W$}
}
\caption{Verification of a certificate $(\Tbl,\Lv,\Rmax)$ by exhaustive expansion from the root.}
\label{alg:verify}
\end{algorithm}

Some comments on \cref{alg:verify}.

\begin{itemize}[leftmargin=1.6em]
\item \emph{The children.} Lines \ref{ln:childa}--\ref{ln:childd} are a literal transcription of
  \cref{def:states}. In a danger child the
  four revealed stones on the line $a$ form a set $Q_a\subseteq a\setminus\{c\}$ with $F\cap a\subseteq Q_a$;
  since $a\setminus\{c\}$ has five cells, $Q_a$ is determined by the one cell $p$ that is left out, and the
  condition $F\cap a\subseteq Q_a$ says $p\notin F$; likewise for $Q_b$ and $q$. All lines are treated alike, all pairs of live lines
  through a common cell are used, and no child is omitted for reasons of symmetry or cost; symmetric
  children are merged later, by the set $\mathit{Seen}$. The stone bounds admit shortcuts that are
  identities, not heuristics: for $e\notin F$ one has $|F\cup(l\setminus\{e\})|=|F|+5-|F\cap l|$, and
  $|F\cup Q_a\cup Q_b|=|F|+(4-|F\cap a|)+(4-|F\cap b|)$ because $Q_a\cap Q_b=\emptyset$.
\item \emph{Shapes.} Every generated state has the right shape ($|B|=t$ for P-states, $|B|=t-1$ for
  Q-states, $|F|\le t$, $B\cap F=\emptyset$); this follows from the conditions under which the children are
  generated. The implementations assert it.
\item \emph{The \VOID{} test.} The only place where a generated child is discarded is line~\ref{ln:void}, and the
  test uses the child's own sets. The child (a) and the child (i) are never \VOID{} when their
  parent is not: they have the same set $F$ as the parent and at least the stones $B$ of the parent
  (\cref{lem:adm-closed}). So line~\ref{ln:void} never discards such a child, in agreement with rule (R2)
  of \cref{def:certificate}, which puts the child (i) into $V$ without a condition. An implementation
  should treat a \VOID{} child of type (a) or (i) as an internal error and reject; the implementations of
  \cref{tab:verifiers} that test for it do so, and the case never occurred.
\item \emph{Counting.} The numbers reported in \cref{tab:tree} are obtained as follows. States are
  counted when they are added to $\mathit{Seen}$, at their round $t$. A child of type (c) or (d) is
  counted, at the round $t+1$ of the child, every time \textsc{Push} is called for it in lines
  \ref{ln:childa}--\ref{ln:childd} and does not discard it; by \cref{lem:children-distinct} these calls
  are made for pairwise different children of the state being expanded, so no child is counted twice for
  the same parent, while children of different states with the same canonical form are counted
  separately. A \VOID{} child of type (ii) is counted at the round $t$ of its parent.
\item \emph{The table move} $\Tbl(s)$ is a cell in the canonical frame of $s$. As $W$ contains canonical
  states only, it is applied directly; no move is ever translated to another frame.
\item \emph{$\LOST$.} For $\Rmax=10$ the requirement in line~\ref{ln:lost} is vacuous: an expanded state has
  $t\le 9$, and by \cref{prop:lost-stones} no pattern $\Etwo$ or $\Ethree$ fits into $t+1\le 10$ stones. The two
  verifiers of \cref{ssec:impl} implement the test of \cref{def:lost} nevertheless (a cheap count of
  missing stones rejects almost all states at once), so that they can be used for larger $\Rmax$.
\item \emph{Termination.} A P-state with $t\ge\Rmax$ is never expanded, so only finitely many states
  occur; each is processed once.
\end{itemize}

\begin{proposition}[Correctness of the verifier]\label{prop:verify-correct}
Suppose that \cref{alg:verify}, carried out exactly on the input $(\Tbl,\Lv,\Rmax)$, returns \accept.
Then $(\Tbl,\Lv,\Rmax)$ is a valid certificate, and hence the root is \WON{} for this $\Rmax$.
\end{proposition}

\begin{proof}
Let $V$ be the reachable set. Every state that is ever put into $\mathit{Seen}$ belongs to $V$: this holds
for the root, and a state is added only as the canonical form of a non-\VOID{} child of a state of
$\mathit{Seen}$ that is an unlisted P-state or a Q-state with a legal table move, exactly as in the rules
(R1) and (R2). Conversely, at the end every state of $\mathit{Seen}$ has been removed from $W$ and
processed without a failed requirement, so $\mathit{Seen}$ contains the root and is closed under (R1) and
(R2); by minimality $V\subseteq\mathit{Seen}$. Hence $\mathit{Seen}=V$, and the requirements that were
checked when the states were processed are (V1), (V2) and (V3).
\end{proof}

\begin{remark}[Representatives that are not minimal]\label{rem:any-image}
Suppose that in \cref{alg:verify} the function $\canon$ is replaced by any function $r$ that maps every
state to the image of this state under some element of $\G$, which may depend on the state, and that
the algorithm returns \accept. Then the root is still \WON. Indeed, every state that was processed is
the root or $r(s')$ for a child $s'$ that is not \VOID, hence it is not \VOID{}
(\cref{lem:game-invariance}\,(a)); and every such state $s$
is \WON, by induction on the rank of $s$ (\cref{lem:rank}; a state and its images under $\G$ have the
same rank, since the rank depends only on $t$, $|F|$ and the kind of the state). If $s$ is a P-state
in $\Lv$, then $\LEAF$ was required. If $s$ is a P-state outside $\Lv$, then $t<\Rmax$ and the
falsity of $\LOST$ were required, and for every child $s'$ of $s$ of type (a), (c) or (d) that is not
\VOID{} the state $r(s')$ was processed and has smaller rank than $s$; so $r(s')$ is \WON, and so is
$s'$ by \cref{lem:game-invariance}\,(c). If $s$ is a Q-state, the cell $m=\Tbl(s)\notin B\cup F$ that
was required is a table move by the same argument, applied to the child (i) and to the child (ii) if
it exists and is not \VOID. Finally the root is processed, because it is fixed by every element of
$\G$. In this situation the triple $(\Tbl,\Lv,\Rmax)$ need not be a valid certificate in the sense of
\cref{def:certificate}, whose reachable set consists of canonical forms; minimality of the stored
representative matters for finding the entries of the two files, which are listed by canonical form.
\end{remark}

\subsection{The leaf test}\label{ssec:algleaf}

Recall from \cref{def:leaf} that $\LEAF(B,F,k)$ is the conjunction of (L1) $F$ is $B$-admissible, (L2)
a pairing for $(B,F)$ exists, and (L3) every $F$-connected set $S$ of at least three live lines has
$\sur_F(S)\ge\theta$, where $\theta=k-|F|$. Condition (L3) quantifies over exponentially many sets. By
\cref{prop:triples} it is equivalent, given (L2), to $\mu(T)\ge\theta$ for every $F$-connected triple $T$
of live lines, where $\mu(T)$ is the minimum of $\sur_F(S)$ over all sets $S\supseteq T$ of live lines; and
by \cref{prop:maxflow} $\mu(T)$ is the number of augmentations that can be performed, starting from a
pairing, when the three lines of $T$ are allowed to hold arbitrarily many cells. \Cref{alg:leaf} spells
out the resulting test. For
$\theta\ge 1$ a line $l$ is called \emph{strong} if $\mu(\{l\})\ge\theta$, and a pair $\{a,b\}$ of lines is
called strong if $\mu(\{a,b\})\ge\theta$; a triple that contains a strong line or a strong pair has
$\mu(T)\ge\theta$ (\cref{cor:pruning}) and need not be tested.

We use the notions of \cref{def:U-assignment}. The algorithm maintains an assignment $\mathit{own}$ of cells
to live lines: $\mathit{own}(c)=l$ means that the line $l$ holds the cell $c\in l\setminus F$, and every
cell is held by at most one line. An augmenting path from a line $l_0$ consists of distinct lines
$l_0,\dots,l_j$ and cells $c_1,\dots,c_{j+1}$ with $c_i\in l_{i-1}\setminus F$ for all $i$,
$\mathit{own}(c_i)=l_i$ for $i\le j$, and $c_{j+1}$ not assigned; augmenting along it sets
$\mathit{own}(c_i)\gets l_{i-1}$ for $i=1,\dots,j+1$, so that $l_0$ gains one cell and every other line holds
as many cells as before. An augmenting path \emph{from a set $U$} of lines is one from some line of $U$.
Such a path is found by a breadth-first search over the lines.

\begin{algorithm}[t]
\DontPrintSemicolon
\SetKwProg{Fn}{function}{}{}
\SetKwFunction{Gain}{Gain}
\SetKwFunction{Aug}{Augment}
\KwIn{disjoint sets $B$, $F$ of cells and an integer $k\ge|F|$}
\KwOut{the truth value of $\LEAF(B,F,k)$}
\lIf{$F$ is not $B$-admissible}{\Return false}\nllabel{ln:adm}
$L\gets$ the set of live lines;\quad $\mathit{own}\gets$ the empty assignment\;
\ForEach{$l\in L$}{\nllabel{ln:pair1}
  \For(\tcp*[h]{every live line needs two cells}){$i\gets 1$ \KwTo $2$}{
    \lIf{\Aug{$\mathit{own},\{l\}$} fails}{\Return false}\nllabel{ln:pair2}
  }
}
\tcp{now $\mathit{own}$ is a pairing for $(B,F)$: every live line holds exactly two cells}
$\theta\gets k-|F|$\;
\lIf{$\theta\le 0$}{\Return true}
$L_w\gets\{\,l\in L:\ \Gain{$\mathit{own},\{l\},\theta$}<\theta\,\}$ \tcp*[r]{the lines that are not strong}\nllabel{ln:strong}
\ForEach{triple $T=\{l_1,l_2,l_3\}\subseteq L_w$ in which $l_2$ shares a cell outside $F$ with $l_1$ and with $l_3$}{\nllabel{ln:triples}
  \lIf{\Gain{$\mathit{own},T,\theta$} $<\theta$}{\Return false}
}
\Return true\;
\BlankLine
\Fn{\Gain{$\mathit{own},U,\theta$}}{
  $\mathit{own}'\gets$ a copy of $\mathit{own}$;\quad $n\gets 0$\;
  \lWhile{$n<\theta$ and \Aug{$\mathit{own}',U$} succeeds}{$n\gets n+1$}
  \Return $n$\;
}
\Fn{\Aug{$\mathit{own},U$}}{
  search for an augmenting path from $U$ with respect to $\mathit{own}$\;
  \leIf{one is found}{augment $\mathit{own}$ along it and \Return success}{\Return failure}
}
\caption{The test $\LEAF(B,F,k)$ by integer max-flow.}
\label{alg:leaf}
\end{algorithm}

\needlines{8}
\begin{proposition}[Correctness of the leaf test]\label{prop:leaf-correct}
\Cref{alg:leaf} returns true if and only if $\LEAF(B,F,k)$ holds.
\end{proposition}
\nopagebreak

\begin{proof}
Line~\ref{ln:adm} is condition (L1). The loop of lines \ref{ln:pair1}--\ref{ln:pair2} decides (L2). Throughout the loop every cell is held
by at most one line, a line holds only cells of its own outside $F$, and every line holds at most two
cells: an augmentation from $l$ gives one more cell to $l$, which held at most one, and leaves the
number of cells of every other line unchanged. If the loop ends, every live line holds exactly two
cells, and $\mathit{own}$ is a pairing for $(B,F)$. If \textsc{Augment}$(\mathit{own},\{l\})$ fails in
line~\ref{ln:pair2}, let $R$ be the set of lines that are reachable from $l$ by alternating paths
(\cref{def:U-assignment}; the definition of alternating paths and fact~(iii) after it do not require
that the lines hold two cells). Every cell of $\N(R)\setminus F$ is held by a line of $R$: a cell
$c\in l'\setminus F$ with $l'\in R$ that is not assigned would complete an augmenting path from $l$,
and a cell that is held by a line $l''\ne l'$ makes $l''$ reachable. At this moment every line of $R$
holds at most two cells and $l\in R$ holds at most one; hence $|\N(R)\setminus F|\le 2|R|-1$, and there
is no pairing for $(B,F)$ by \cref{lem:hall}\,(b). If $\theta\le 0$, condition (L3)
follows from the pairing, since then $\sur_F(S)\ge 0$ for every $S$ (\cref{lem:hall}). Let $\theta\ge1$. By
\cref{prop:maxflow}\,(d), $\textsc{Gain}(\mathit{own},U,\theta)=\min(\theta,\mu(U))$ for every set $U$ of
live lines, independently of the order in which augmenting paths are chosen. A set of three live lines is
$F$-connected if and only if one of them is $F$-adjacent to the other two, so
line~\ref{ln:triples} enumerates exactly the $F$-connected triples without a strong line
(line~\ref{ln:strong}). A triple that contains a strong line has
$\mu(T)\ge\theta$ by \cref{cor:pruning}\,(a). Hence the algorithm returns true if and only if
$\mu(T)\ge\theta$ for all $F$-connected triples $T$, which by \cref{prop:triples}\,(b) is condition (L3).
\end{proof}

\begin{remark}[True verdicts carry explicit witnesses]\label{rem:true-cert}
For the soundness of a verifier only the implication ``the algorithm returns true $\Rightarrow\LEAF$''
matters, and this implication does not need the max-flow min-cut theorem. Suppose
$\textsc{Gain}(\mathit{own},U,\theta)=\theta$, and let $\mathit{own}'$ be the final assignment: no cell is
assigned twice, every live line outside $U$ holds exactly two cells, and the lines of $U$ hold $2|U|+\theta$
cells together. For every set $S\supseteq U$ of live lines, the lines of $S$ then hold $2|S|+\theta$
distinct cells of $\N(S)\setminus F$, so $\sur_F(S)\ge\theta$ (\cref{prop:maxflow}\,(a)). Every $F$-connected
set $S$ with $|S|\ge3$ contains an $F$-connected triple $T$ (\cref{lem:triple}), and $T$ either has
$\textsc{Gain}=\theta$ itself, or contains a strong line $l$, or, if the second pruning rule below is
used, contains a pair $\{a,b\}$ with $\textsc{Gain}=\theta$; in the last two cases the same counting
applies with $U=\{l\}$ or $U=\{a,b\}$. An implementation can therefore validate each true verdict
literally, by checking the final assignments against these three properties. The main verifier does
so, in the modified form described under ``Deleting strong lines'' below (\cref{lem:deleted}). False
verdicts affect only completeness: a wrong ``false'' can make a verifier reject a valid certificate,
never accept an invalid one.
\end{remark}

\subsubsection*{Accelerations}
\Cref{cor:pruning} permits a second pruning rule: a triple containing two lines $a,b$ with
$\mu(\{a,b\})\ge\theta$ need not be tested. A pairwise $F$-adjacent triple is found three times in
line~\ref{ln:triples} and needs to be tested once. The copy in \textsc{Gain} can be avoided by undoing
the augmentations. None of this is needed for correctness, and independent implementations differ in what they use
(\cref{tab:verifiers}).

\subsubsection*{Deleting strong lines}
The main verifier uses one further device, and this one changes the assignments on which its true
verdicts rest. Let $\theta\ge1$. As soon as a line has been found strong it is \emph{deleted}: it gives
up its two cells and is ignored in all later searches, for single lines, pairs and triples alike. To
justify this, let $D\subseteq\live{B}$ be a set of lines and $L_D=\live{B}\setminus D$. For
$T\subseteq L_D$, a \emph{$T$-assignment relative to $D$} is a map $\alpha$ from a set of cells to $L_D$
such that $c\in\alpha(c)$ and $c\notin F$ for every assigned cell $c$, every line of $L_D\setminus T$
holds exactly two cells and every line of $T$ holds at least two; its \emph{excess} is
$\sum_{l\in T}(|\alpha^{-1}(l)|-2)$. The lines of $D$ hold no cell. Put
\[
  \mu_D(T)=\min\{\sur_F(S) : T\subseteq S\subseteq L_D\}.
\]
For $D=\emptyset$ these are the notions of \cref{def:U-assignment} and the number $\mu(T)$.

\begin{lemma}[Deleted lines]\label{lem:deleted}
Let $\theta\ge 1$, let $D\subseteq\live{B}$ be a set of lines with $\mu(\{d\})\ge\theta$ for every
$d\in D$, and let $T\subseteq L_D$.
\begin{enumerate}[label=\textup{(\alph*)}]
\item If there is a $T$-assignment relative to $D$ of excess at least $\theta$, then $\mu(T)\ge\theta$.
\item $\min(\theta,\mu(T))=\min(\theta,\mu_D(T))$. If $\pi$ is a pairing for $(B,F)$, then $\mu_D(T)$ is
      the number of augmentations along augmenting paths from lines of $T$ that are performed when one
      starts from the restriction of $\pi$ to the cells held by the lines of $L_D$, uses lines of $L_D$
      only, and continues until there is no such path.
\end{enumerate}
\end{lemma}

\begin{proof}
(a) Let $T\subseteq S\subseteq\live{B}$; we have to show $\sur_F(S)\ge\theta$. If $S$ contains a line
$d\in D$, then $\sur_F(S)\ge\mu(\{d\})\ge\theta$ by the definition of $\mu$. Otherwise $S\subseteq L_D$.
The cells held by the lines of $S$ are distinct cells of $\N(S)\setminus F$, and their number is at least
$2|S|+\theta$, because $T\subseteq S$, the lines of $S\setminus T$ hold two cells each and the lines of
$T$ hold at least $2|T|+\theta$ cells together. Hence $|\N(S)\setminus F|\ge 2|S|+\theta$.

(b) By definition, $\mu(T)$ is the smaller of $\mu_D(T)$ and of the minimum of $\sur_F(S)$ over the sets
$S\supseteq T$ of live lines that contain a line of $D$; the latter is at least $\theta$, as in (a). This
gives the first assertion. \Cref{def:U-assignment}, \cref{prop:maxflow} and its proof use no property of
the set $\live{B}$ except that $\pi$ assigns two cells of $l\setminus F$ to each of its lines $l$, all
these cells being distinct; they apply verbatim to the set $L_D$ of lines and to the restriction of
$\pi$. Part (d) of the proposition, read for $L_D$, is the second assertion.
\end{proof}

By induction on the order of deletion, every deleted line is indeed strong: when a line $d$ is tested,
the lines deleted before it are strong by the induction hypothesis, and if the test finds a
$\{d\}$-assignment of excess $\theta$ relative to them, then $\mu(\{d\})\ge\theta$ by part (a). Hence
\cref{lem:deleted} applies to every set $D$ of deleted lines that occurs, and single lines, pairs and
triples may be tested in the network without the deleted lines: by (a) a true verdict of such a test is
correct, and by (b) so is a false one. For a verifier that deletes strong lines the properties to be
validated behind a true verdict for a set $T$ are therefore: no cell is assigned twice; every assigned
cell lies on its line and outside $F$; the deleted lines hold no cell; every other live line outside $T$
holds exactly two cells; every line of $T$ holds at least two cells; and the lines of $T$ hold
$2|T|+\theta$ cells together.

\subsubsection*{The leaf test of the main verifier, and its cost}
What \fn{verify/treecheck.cpp} does after it has checked (L1) and computed a pairing is stated step by
step in \cref{ssec:supp-mainleaf}, together with the numbers of flows that it computes. In brief: an
$F$-connected component of at least three live lines whose surplus is less than $\theta$ leads to the
verdict false at once (it violates (L3) of \cref{def:leaf} by definition); a line $l$ with $|l\setminus F|-2<\theta$ is treated as
not strong without a flow; for $\theta=1$ the strong lines are found by one backward search from the
unassigned cells; every line found strong is deleted; and flows are then computed for pairs and for
triples of the remaining lines. All steps are justified by \cref{cor:pruning}, \cref{rem:true-cert} and
\cref{lem:deleted}, and every flow that is used for a true verdict is validated literally, in the form
relative to the deleted lines stated above. \fn{verify/spotcheck.py} implements \cref{alg:leaf} with both
pruning rules and without deletion. The cost of the test is small (\cref{ssec:supp-mainleaf}): in our
runs a leaf test takes about
$0.1$ milliseconds in the C++ verifier (with the literal validation of all flows) and about $2$
milliseconds in the Python verifier.

\subsection{Independent implementations and their results}\label{ssec:impl}

\Cref{tab:verifiers} lists the programs that have verified the certificate. The first two are the
verifiers of the archive and are the ones a reader should run; they share no code with each other, with
the search, or with the simulator. All programs of the table are contained in the archive
(\cref{tab:artifact}).

{\footnotesize
\begin{longtable}{@{}>{\raggedright\arraybackslash}p{0.215\textwidth}>{\raggedright\arraybackslash}p{0.435\textwidth}>{\raggedright\arraybackslash}p{0.29\textwidth}@{}}
\toprule
program & what it does & result\\
\midrule
\endfirsthead
\toprule
program & what it does & result\\
\midrule
\endhead
\midrule
\multicolumn{3}{@{}r@{}}{\emph{continued on the next page}}\\
\endfoot
\bottomrule
\addlinespace[6pt]
\caption{Programs that have verified the certificate, and their results. ``Same numbers in every round''
means: the numbers of P-states, leaves and Q-states in every round are those of \cref{tab:tree}; ``for
every type of child'' means that this also holds for the threat children, the danger children and the
\VOID{} children. Sources of the numbers: \fn{results/final/treecheck.json},
\fn{treecheck_crosscheck.json}, \fn{spotcheck_full.json}, \fn{stats.json}; the review reports
\fn{logs/review/treecheck_audit.md} and \fn{REVIEW.md} with the outputs \fn{aexpand_final.json} and
\fn{lead_expand.json}; and, in the directory \fn{logs/paper/}, the reports \fn{referee_replication.md},
\fn{referee_audit.md} and \fn{referee_math_b.md} of the internal refereeing.}
\label{tab:verifiers}\\
\endlastfoot
\multicolumn{3}{@{}l}{\emph{The two verifiers of the archive}}\\
\addlinespace[2pt]
\fn{verify/treecheck.cpp} ($4{,}012$ lines of C++), option \fn{--exact}
 & \Cref{alg:verify} with its own geometry, group, canonical form, max-flow and $\LOST$; visited set keyed
   by full states; \cref{alg:leaf} with both pruning rules and deletion of strong lines
   (\cref{lem:deleted}); for $\theta=1$ the strong lines are found by one backward search; every
   assignment behind a true verdict validated literally (\cref{rem:true-cert})
 & \accept; $246{,}299$ states, $106{,}723$ leaves recomputed, $0$ errors; $8.2$ CPU-seconds in the
   run of \fn{treecheck.json}; reference timing in \cref{ssec:certifying}\\
\addlinespace
\quad the same with \fn{--crosscheck}
 & in addition: the canonical form of every state by brute force over the $1{,}152$ maps; children and $\LOST$
   of all $46{,}412$ expanded P-states against literal reference code; $\LEAF$ of $2{,}153$ leaves by an
   unpruned enumeration of triples with separate matching code
 & \accept; same numbers in every round; $945$\,s on two threads\\
\addlinespace
\fn{verify/spotcheck.py} ($1{,}496$ lines of Python), exhaustive mode
 & \Cref{alg:verify} and \cref{alg:leaf} (pruning by lines and pairs) written from the specification
   alone; exact keys everywhere; also checks that every entry of both files is canonical and is reached
 & \accept; same states and \VOID{} children in every round; children (a)/(c)/(d)/(i)/(ii) in total:
   $46{,}412$ / $143{,}498$ / $28{,}275$ / $93{,}164$ / $54{,}482$, of which $151$ \VOID; $339$
   CPU-seconds\\
\midrule
\multicolumn{3}{@{}l}{\emph{Written during the internal review of the project}}\\
\addlinespace[2pt]
re-expansion I (\fn{aexpand.cpp}, \fn{aleaf.cpp}; $450+227$ lines of C++)
 & expansion with brute-force canonical forms, exact keys and literal children; separate program for
   $\LEAF$ with an explicit flow network, all triples without pruning for $\theta\ge2$
 & same numbers in every round and for every type of child; all $106{,}723$ leaves satisfy $\LEAF$;
   about $5$ CPU-minutes\\
\addlinespace
re-expansion II (\fn{lead_expand.cpp}, \fn{lead_leaf.cpp}; $367+307$ lines of C++)
 & expansion with brute-force canonical forms, exact keys and literal children ($\LOST$ not evaluated; the
   program refuses $\Rmax>10$); separate program for $\LEAF$, pruning by strong lines only, and no pruning
   at all on the $1{,}912$ leaves with $\theta\ge3$
 & $0$ errors; same numbers in every round and for every type of child; all leaves satisfy $\LEAF$;
   $44+44$ CPU-seconds\\
\midrule
\multicolumn{3}{@{}l}{\emph{Written from the text of this paper during the internal refereeing of its first version}}\\
\addlinespace[2pt]
replication (\fn{verify.cpp}; $706$ lines of C++)
 & written by a referee who was given only the PDF of the paper and the two certificate files:
   \cref{alg:verify,alg:leaf} with brute-force canonical forms, exact keys, literal children, a literal
   $\LOST$ at all expanded states, and all assignments behind true verdicts checked as in
   \cref{rem:true-cert}
 & \accept; $246{,}299$ states; every entry of \cref{tab:tree,tab:leaf-theta} reproduced; $42$
   CPU-seconds on one thread (exit code $1$ for the reason given in \cref{ssec:own})\\
\addlinespace
audit (\fn{ref_audit.cpp}; $474$ lines of C++)
 & written by the referee who audited the numbers of this paper, from
   \cref{sec:board,sec:leaf,sec:game}, \cref{alg:verify,alg:leaf} and \cref{sec:formats}:
   brute-force canonical forms, exact keys, literal children, \cref{alg:leaf} with pruning by strong
   lines only on every leaf ($\LOST$ not evaluated; the program refuses to expand a state with $t>9$)
 & \accept, $0$ errors; $246{,}299$ states; every entry of \cref{tab:tree} reproduced; $63$
   CPU-seconds\\
\addlinespace
transcription of \cref{alg:verify} (\fn{b_expand.py}; $161$ lines of Python and two modules)
 & written by a third referee from the printed pseudocode and \cref{def:states,def:certificate}, with
   its own geometry and canonical form; \emph{no leaf test}: of $\LEAF$ only (L1) is checked on all
   leaves and (L2) on $1{,}500$ of them
 & every number of \cref{ssec:fmt-targets,ssec:supp-targets} reproduced; every entry of both files
   reached; $3.4$ minutes\\
\midrule
\multicolumn{3}{@{}l}{\emph{Not independent of the search}}\\
\addlinespace[2pt]
\fn{bin/search} with \fn{--selfcheck}
 & re-walk of the two files by the search program with a naive child generator and a literal $\LOST$;
   shares geometry, canonical form and $\LEAF$ code with the search
 & $46{,}412$ expanded P-states, $93{,}164$ Q-states, $106{,}723$ leaves, $0$ errors (totals only)\\
\end{longtable}}

\subsubsection*{Independence}
Throughout this paper a program is called \emph{independently written} if it was written in a separate
session, from the written specification or from the text of this paper, without access to the source
code of the other programs. The declaration ``Use of AI-based tools'' at the end of the paper says how
the programs were produced: all of them were written by AI-based software agents, and the internal
review of the project and the internal refereeing of the first version of this paper, to which this
section refers, were carried out by such agents as well. A ``referee'' is, in this section, in
\cref{sec:repro,sec:formats} and in the supplementary material, a session of the internal refereeing.

In this sense the two verifiers of the archive are independent of each other, of the search and of the
simulator. The first, \fn{verify/treecheck.cpp}, is a single self-contained file; the second,
\fn{verify/spotcheck.py}, was written
from the written specification alone (the definitions reproduced in
\cref{sec:board,sec:pairing,sec:leaf,sec:game} and \cref{sec:formats}). The two re-expansions were written
during the internal review, in the sessions that audited the main verifier and that reviewed the
project as a whole. These sessions had read programs under review, so the re-expansions are independent
only in the weaker sense that they share no code with any other program: each has its own geometry and
computes canonical forms by brute force over all $1{,}152$ maps. They were written quickly and were
tested mainly against the other programs. The verifier listed as ``replication'' was written during the
internal refereeing from the PDF of the paper (and a text extraction of it) and the two certificate
files alone: its author had neither the specification nor the sources of the paper nor any program. The programs listed as ``audit'' and ``transcription'' were written from the text of
the paper as well, by referees who had the project files at their disposal; they share no code with
the programs of the project. The self-check of the search is independent of all the others but
not of the search, with which it shares its library. All these programs use integer arithmetic only.

Independence in this sense does not exclude a common misreading of the specification, the less so as all
programs were produced with tools of the same kind. Against that stand the precise statement of the
algorithms in this paper, the agreement of the numbers in every round, the comparisons with brute force
of \cref{ssec:audits}, the verifier that was written from the paper alone, and the possibility of writing
one's own verifier (\cref{ssec:own}).

\subsubsection*{Agreement}
All programs accept. The two \fn{treecheck} runs, the two re-expansions and the replication and audit
verifiers report the same numbers in every round and for every type of child (\cref{tab:tree});
\fn{spotcheck.py} reports the same states and \VOID{} children in every round and the same totals of
children; the transcription of \cref{alg:verify} reports the same states, threat children, danger
children and \VOID{} children in every round; the self-check of the search reports the same totals. Every
one of the $93{,}164$ table entries and $106{,}723$ leaf entries is canonical under seven separately
written implementations of the canonical form (all programs of \cref{tab:verifiers} except the last), is
reached, and is used. $\LEAF$ has been recomputed for every leaf by seven implementations (the two
verifiers, the two re-expansions, the replication and audit verifiers, and the library of the search),
six of which share no code with the search.

When a result of \fn{treecheck} is cited, the following fields of its output should be checked in addition
to \texttt{ok}: the root is \texttt{"0 - -"}, \texttt{root\_void} is false, \texttt{complete} is true,
\texttt{rmax} is $10$, \texttt{exact\_visited\_set} and \texttt{paranoid\_flows} are true, the numbers of
non-canonical entries, of malformed entries and of table conflicts are $0$, and the byte counts and checksums of the
input files are the expected ones. All of this holds for \fn{results/final/treecheck.json}.

\subsection{Audits of the main verifier}\label{ssec:audits}

Since \fn{verify/treecheck.cpp} is the largest of the verifiers, it was audited separately
(\fn{logs/review/treecheck_audit.md}). The audit read the code against the eight dangerous defects listed in
\cref{ssec:what} and made three kinds of experiments, which are reported in detail in
\cref{ssec:supp-audits}.

\emph{Mutations of the certificate.} In all, $66$ seeded mutations of the two files or of the options were
submitted to the verifier. Of these, $58$ were rejected, among them every deleted table entry, every
replaced table move, every deleted leaf, every bogus leaf and the bounds $\Rmax=9$ and $\Rmax=8$. Eight
were accepted. None of the eight removes or changes an obligation: in each of them the visited states and
all counters are those of the unmodified run.

\emph{Mutations of the verifier.} Twelve variants of the source with one seeded defect each were compiled
and run. The self-test of the verifier detects all $12$. On the real certificate one mutant changes the
verdict and seven change the counters of \cref{tab:tree}; four change neither: a constant-false $\LOST$,
which is harmless for $\Rmax=10$, and three defects that concern the leaf test. This last finding is the
reason why the correctness of the leaf test is not inferred from the run of one program: a valid
certificate does not exercise a wrong ``true'', so $\LEAF$ is covered by independent implementations and by
comparisons with brute force instead.

\emph{Comparisons with brute force.} Three leaf tests were compared with a direct evaluation of
\cref{def:leaf} over \emph{all} sets of live lines, on instances with few live lines: that of
\fn{treecheck} on $53{,}755$ instances, that of \fn{spotcheck.py} on $6{,}000$ and that of the replication
verifier of \cref{tab:verifiers} on $19{,}995$. The leaf test of \fn{treecheck} was compared with that
of the library of the search and with that of the audit verifier of \cref{tab:verifiers}, in verdict and
in the failing condition, on all leaves and on $20{,}000$ perturbed instances derived from the leaves. The
predicate $\LOST$ of \fn{treecheck} was compared with a literal transcription of the patterns $\Etwo$ and
$\Ethree$ on $800$ random admissible states. No comparison found a mismatch.

\subsection{Cross-checks that are not part of the proof}\label{ssec:cross}

The following computations test the certificate and the programs from other directions. None of them is
needed for \cref{thm:main}. The programs, the result files and the exact coverage are given in
\cref{ssec:supp-cross}; \cref{tab:leaves} shows how each leaf is covered.

\begin{table}[ht]
\centering\small
\begin{tabular}{@{}rrrrr@{}}
\toprule
 & & $\LEAF$ by exact & $\CC$ by exact & $\CC$ by the integer\\
$\theta$ & leaves & max-flow & enumeration & program (floating point)\\
\midrule
0 & 64{,}841 & all & all & 300 sampled\\
1 & 32{,}703 & all & all & 300 sampled\\
2 & 7{,}267 & all & all & 150 sampled\\
3 & 1{,}434 & all & all & 80 sampled\\
4 & 343 & all & --- & all\\
5 & 105 & all & --- & all\\
6 & 25 & all & --- & all\\
7 & 4 & all & --- & all\\
8 & 1 & all & --- & all\\
\midrule
total & 106{,}723 & 106{,}723 & 106{,}245 & 1{,}308\\
\bottomrule
\end{tabular}
\caption{Coverage of the leaves by $\theta=t-|F|$. The third column is the proof obligation (V1),
established by every program of \cref{tab:verifiers} that contains a leaf test; the last two columns are
cross-checks of the conclusion $\CC(B,F,t)$ of \cref{thm:leaf-implies-C}. All answers are positive.}
\label{tab:leaves}
\end{table}

\emph{Exact enumeration of the pairing condition.} $\LEAF$ is only sufficient for the pairing condition
$\CC(B,F,k)$ (\cref{thm:leaf-implies-C}). For small $\theta$ the statement $\CC(B,F,t)$ itself can be
decided by enumerating the admissible sets (\cref{rem:C}\,(iv)). This was done for all leaves with
$\theta\le3$, that is, for $106{,}245$ of the $106{,}723$ leaves ($99.55\%$): the pairing condition holds
for every one of them.

\emph{The integer program.} \Cref{sec:ip} describes a $0$--$1$ program whose optimal value is positive if
and only if $\CC(B,F,k)$ fails. It was solved with the floating-point solver
HiGHS~\cite{HuangfuHall2018} on all $478$ leaves with $\theta\ge4$ and on a seeded random sample of $830$
leaves with $\theta\le 3$: all $1{,}308$ answers are ``holds''. On $1{,}291$ random and structured
instances it was compared with an exact integer decision procedure for $\CC$; wherever both gave an
answer, the answers agree. Among the $656$ of these instances for which $\CC$ holds there are $160$ for
which $\LEAF$ is false: $\LEAF$ is a strictly stronger condition than $\CC$ (it fails on instances on
which $\CC$ holds), not only in the situation of \cref{ex:C-not-leaf}. No such state is needed by the
certificate.

\emph{Simulated games.} The program \fn{sim/sim.cpp}, which was written from the specification and shares
no code with the other programs, plays the real strategy $\sigma$ of \cref{def:sigma} on the real board.
In $10^6$ completed games against a mixture of nine adversaries there was no line of \Maker, no missing
table entry, no leaf whose real position has no pairing, no position that is not admissible after
\Breaker's move and no illegal table move. This is a sampling test and nothing more: the games met
$84.8\%$ of the table entries and $76.7\%$ of the leaves, and the other states are covered by the
verifiers of \cref{ssec:impl} only.


\subsection{What is trusted}\label{ssec:trust}

\Cref{thm:main} is a computer-assisted theorem. Its proof consists of the written proofs of the
preceding sections and of \cref{fact:H}. To accept it, a reader has to trust
\begin{itemize}[leftmargin=1.6em]
\item the written proofs, in particular the soundness theorem (\cref{thm:soundness}) and
  \cref{thm:leaf-implies-C};
\item that the definition of validity used by the verifiers (\cref{ssec:what}) is the one of
  \cref{def:certificate}, and that \cref{alg:verify,alg:leaf} are correct
  (\cref{prop:verify-correct,prop:leaf-correct}), as is the variant of \cref{alg:leaf} with deleted
  lines that the main verifier uses (\cref{lem:deleted});
\item at least one implementation of these algorithms (\fn{verify/spotcheck.py} implements
  \cref{alg:leaf} without deletion of lines, \fn{verify/treecheck.cpp} the variant with deleted lines),
  together with its compiler or interpreter
  (g++~13.3 for the C++ programs, CPython~3.13 for \fn{spotcheck.py}), the operating system and the
  hardware on which it ran. All runs of the project that are reported here were made on one machine; the
  runs of the three programs written during the internal refereeing (\cref{tab:verifiers}) are documented
  in the referee reports.
\end{itemize}
Nothing has been formally verified in a proof assistant. No floating-point computation and no random
choice enters the proof; the integer program and the simulator are cross-checks. All programs, the
verifiers included, and the internal review were produced by AI-based software agents; what
``independent'' means under these circumstances is stated in \cref{ssec:impl} and in the Declarations at
the end of the paper.

\subsubsection*{Hashing}
\fn{spotcheck.py}, the two re-expansions and the three programs written during the internal refereeing
use exact keys throughout. \fn{treecheck} with the option
\fn{--exact} keys its visited set by the full state; without that option it would use $128$-bit hashes, and
a collision could make it skip a state. Its look-ups in the two files always go through $120$-bit hashes.
A collision there cannot make it accept a certificate for which the root is not \WON: a state that is
mistaken for a listed leaf must still pass the leaf test for its own sets $B$, $F$, so it is \WON{};
and a Q-state that receives the move of another entry still has this move checked for legality and both
children verified, so the move is a correct table move for it. In either case the program has verified a
valid certificate, possibly with a larger leaf set or a modified table. (In the certified run the numbers
of distinct keys, of entries of the files and of visited Q-states and leaves coincide, and the exact-key programs use exactly
the same entries.)

\subsubsection*{Known minor issues}
The audits found no defect that affects a result. Some properties of the tools that should be known to
anyone who reuses them are listed in \cref{ssec:supp-issues} and in the file \fn{README.md} of the
archive; they are the reason for the list of output fields to be checked in \cref{ssec:impl}. All the
counters concerned are zero in the certified run.
